\chapter{Soạn đề cùng Kriky}
\label{chuong:soan-de}

Thầy cô nhắn yêu cầu bằng câu bình thường. Kriky hỏi lại nếu thiếu thông tin, đưa ra một dàn ý, rồi soạn
đề khi thầy cô đồng ý. Đề soạn xong là đề nháp: học sinh chưa thấy gì cho tới khi thầy cô duyệt và phát
hành.

\viec{Nhắn yêu cầu ra đề}

\anh{T/14-chat-moi}{
  \khung{0.01}{0.69}{0.99}{0.81}
  \diem{1}{0.015}{0.69}
  \diem{2}{0.89}{0.89}
}

\begin{buoc}
  \item Bấm \nut{Đoạn chat mới}. Bấm một câu gợi ý để điền sẵn, hoặc tự gõ.
  \item Gõ yêu cầu vào ô \nhan{Nhắn cho Kriky…} rồi bấm \nut{Gửi}.
\end{buoc}

Một yêu cầu tốt nói đủ năm điều:

\begin{loinhan}
Tạo đề kiểm tra 15 phút bài Tính đơn điệu và cực trị của hàm số cho lớp 12A, môn Toán, 8 câu trắc
nghiệm, mức nhận biết và thông hiểu.
\end{loinhan}

\begin{bangtra}{L{3cm}L{6.2cm}Y}{\dau{Nói rõ} & \dau{Ví dụ} & \dau{Thiếu thì}}
  Môn và khối & \go{Toán}, \go{khối 12} & Kriky hỏi lại. Lớp nào nhận đề thì chọn lúc phát hành. \\
  Phạm vi kiến thức & \go{bài Tính đơn điệu và cực trị} & Kriky hỏi lại. \\
  Số câu & \go{8 câu trắc nghiệm} & Kriky hỏi lại. \\
  Mức độ & \go{nhận biết và thông hiểu} & Kriky tự chọn. \\
  Tài liệu & gắn sách vào lời nhắn, hoặc \go{không dùng tài liệu} & Kriky có thể hỏi có dùng tài liệu
  không. \\
\end{bangtra}

\viec{Ra đề theo một tài liệu}

Muốn câu hỏi bám đúng một cuốn sách đã tải lên, thầy cô gắn sách ấy vào lời nhắn.

\anh{T/15-gan-tai-lieu}{
  \diem{2}{0.36}{0.15}
  \diem{3}{0.89}{0.62}
}

\begin{buoc}
  \item Kéo tài liệu (đã đọc xong) từ nhóm \nhan{TÀI LIỆU} thả vào ô nhắn.
  \item Tên tài liệu hiện ngay trên ô nhắn. Gắn nhầm thì bấm \nut{Bỏ}.
  \item Gõ yêu cầu rồi bấm \nut{Gửi}.
\end{buoc}

\viec{Trả lời khi Kriky hỏi lại}

Thiếu thông tin thì Kriky hỏi lại, kèm các nút để chọn nhanh.

\anh{T/16-hoi-lai}{
  \khung{0.03}{0.47}{0.97}{0.95}
  \diem{1}{0.03}{0.47}
}

\begin{buoc}
  \item Bấm nút đúng ý. Không có nút nào đúng thì bấm \nut{Khác…}, gõ ý của mình vào ô \nhan{Ghi rõ
        lựa chọn khác…} rồi bấm \nut{Gửi} (khi Kriky hỏi nhiều câu một lúc, nút này là \nut{Chọn}).
\end{buoc}

Khi Kriky hỏi nhiều câu một lúc, mỗi câu hiện trên một trang của thẻ (\nhan{Câu 1/3}). Chọn xong một câu
thì bấm \nut{Câu sau ›}. Sau câu cuối, bấm \nut{Xem lại ›} để thấy thẻ tóm tắt:

\anh[0.85\linewidth]{T/16b-hoi-lai-tom-tat}{
  \diem{1}{0.88}{0.91}
}

\begin{buoc}
  \item Muốn đổi câu nào thì bấm \nut{Sửa} cạnh câu ấy; đúng hết thì bấm \nut{Gửi}.
\end{buoc}

\viec{Đọc và đồng ý dàn ý}

Trước khi soạn, Kriky đưa ra dàn ý: đề chia mấy phần, mỗi phần mấy câu, hỏi kiến thức gì.

\anh{T/17-dan-y}{}

\begin{buoc}
  \item Đọc dàn ý. Ưng thì \textbf{nhắn} đồng ý, ví dụ \go{Đồng ý, soạn đi.} Dàn ý không có nút đồng ý.
  \item Muốn sửa thì nhắn điều cần đổi, ví dụ \go{Phần 4 đổi thành 2 câu vận dụng}. Kriky đưa dàn ý mới.
\end{buoc}

\xong{Kriky soạn xong, thẻ \nhan{Đã tạo đề "…"} hiện dưới lời nhắn.}

\anh{T/18-da-soan}{
  \diem{1}{0.21}{0.465}
  \diem{2}{0.44}{0.93}
}

Bấm \sobuoc{1} \nhan{Đã làm 2 bước} để xem Kriky đã làm gì. Bấm vào thẻ \sobuoc{2} để mở khung đề.

\machnho{Kriky chưa trả lời được (\nhan{Kriky chưa trả lời được lúc này…}) thì đợi một lát, gõ lại
lời nhắn rồi gửi.}

\viec{Xem đề trong khung đề}

\anh[0.6\linewidth]{T/19-khung-de}{
  \diem{1}{0.73}{0.11}
  \diem{2}{0.755}{0.42}
  \diem{3}{0.33}{0.93}
}

\begin{bangtra}{L{1cm}Y}{\dau{Số} & \dau{Dùng để}}
  \sobuoc{1} & \nut{Đổi tên} đề. Gõ tên mới rồi bấm \nut{Lưu}. \\
  \sobuoc{2} & \nut{Sửa} hoặc \nut{Xoá} một câu. \\
  \sobuoc{3} & \nut{Lời giải · 2 cách}: xem cách giải và lỗi của từng phương án sai. \\
\end{bangtra}

Đáp án đúng có dấu {\color{correct}\faCheck}. Mỗi câu ghi \nhan{Kriky soạn} hay \nhan{Từ ngân hàng}.

\anh[0.8\linewidth]{T/20-loi-giai}{}

Hộp lời giải cho thầy cô đọc từng cách giải, và mỗi phương án sai sai vì đâu (\nhan{MỖI PHƯƠNG ÁN NHIỄU
GẮN MỘT LỖI}). Đọc xong bấm \nut{Đóng}.

\viec{Sửa một câu}

\anh[0.8\linewidth]{T/21-sua-cau}{
  \diem{2}{0.035}{0.71}
  \diem{3}{0.035}{0.79}
  \diem{4}{0.88}{0.8}
}

\begin{buoc}
  \item Bấm \nut{Sửa} ở câu cần sửa. Câu mở thành ô \nhan{Đề bài}.
  \item Mở \nhan{Phương án} để sửa các lựa chọn, bấm nút tròn để chọn đáp án đúng. Mỗi phương án sai cần
        ghi lỗi của nó ở ô \nhan{Lỗi của …}.
  \item Mở \nhan{Cách giải} để sửa lời giải.
  \item Bấm \nut{Lưu}. Không muốn sửa nữa thì bấm \nut{Huỷ}.
\end{buoc}

\neuloi{Dòng đỏ \nhan{Công thức có ký tự hỏng…} như ảnh trên nghĩa là công thức trong câu có chữ lỗi.
Gõ lại công thức ấy, hoặc nhắn Kriky \go{Viết lại công thức câu 2 cho đúng}.}

Mở \nhan{Phương án} thì thấy từng lựa chọn:

\anh[0.55\linewidth]{T/21b-sua-phuong-an}{
  \diem{1}{0.04}{0.15}
  \diem{2}{0.04}{0.4}
  \diem{3}{0.88}{0.92}
}

\begin{bangtra}{L{1cm}Y}{\dau{Số} & \dau{Dùng để}}
  \sobuoc{1} & Nút tròn đánh dấu đáp án đúng. Bấm nút tròn của phương án khác để đổi đáp án. \\
  \sobuoc{2} & Ô \nhan{Lỗi của …} của mỗi phương án sai: ghi vì sao học sinh hay chọn nhầm nó. Các em
  đọc lời này trong phần lời giải. \\
  \sobuoc{3} & \nut{Lưu} khi sửa xong. \\
\end{bangtra}

\machnho{Trong ô sửa, công thức hiện dạng chữ có dấu \go{\$}, ví dụ \go{\$f'(x) > 0\$}. Đó là chữ để sửa,
không phải lỗi; lưu xong nó lại hiện thành công thức.}

\viec{Xoá một câu}

\anh[0.6\linewidth]{T/22-xoa-cau}{
  \diem{1}{0.975}{0.76}
}

Bấm \nut{Xoá} ở câu ấy, rồi bấm \nut{Xoá câu} \sobuoc{1}. Đổi ý thì bấm \nut{Giữ lại}.

\luuy{Câu chưa có trong ngân hàng câu hỏi thì xoá rồi không lấy lại được.}

\viec{Thêm câu từ ngân hàng}

Câu của các đề đã duyệt nằm trong ngân hàng câu hỏi. Thầy cô lấy lại được cho đề mới.

\anh[0.8\linewidth]{T/23-ngan-hang-chon}{
  \diem{1}{0.975}{0.16}
  \diem{2}{0.035}{0.555}
  \diem{3}{0.72}{0.92}
}

\begin{buoc}
  \item Trong khung đề, bấm \nut{\faPlus\ Thêm câu từ ngân hàng}. Gõ vào ô tìm để lọc theo nội dung hay mục
        kiến thức.
  \item Bấm vào các câu muốn lấy. Câu đã chọn có dấu tích; câu \nhan{đã có trong đề này} thì mờ.
  \item Bấm \nut{Thêm 2 câu vào đề} (số câu đổi theo số đã chọn).
\end{buoc}

\viec{Nhờ Kriky soạn thêm hay thay câu}

Nhắn ngay trong đoạn chat của đề, ví dụ:

\begin{loinhan}
Thêm 2 câu mức vận dụng vào đề này.
\end{loinhan}

\begin{loinhan}
Thay câu 5 bằng một câu khác về cực trị.
\end{loinhan}

Kriky đưa dàn ý cho phần soạn thêm, thầy cô đồng ý như trên.

\neuloi{Một câu hiện \nhan{Chưa soạn được}: Kriky bị gián đoạn khi soạn câu ấy. Bấm \nut{Thử lại} ngay
trên câu đó trong khung đề. Các câu khác vẫn sửa, xoá được trong lúc chờ.}

\viec{Đổi tên, xoá đoạn chat}

\anh[0.5\linewidth]{T/24-tuy-chon-doan-chat}{
  \diem{1}{0.955}{0.41}
  \diem{2}{0.955}{0.77}
}

\begin{buoc}
  \item Bấm nút \nut{\faEllipsisH} cạnh một đoạn chat ở thanh trái.
  \item Chọn \nut{Đổi tên}, gõ tên mới tại chỗ rồi nhấn Enter. Hoặc chọn \nut{Xoá} rồi xác nhận bằng
        \nut{Xoá đoạn chat}.
\end{buoc}

\luuy{Đề đã tạo trong đoạn chat vẫn còn khi xoá đoạn chat, nhưng hiện chưa mở lại được từ chỗ khác.
Đừng xoá đoạn chat của một đề chưa phát hành.}

Đề nháp chưa có câu nào thì xoá được bằng \nut{Xoá đề} ở đầu khung đề.
